%!TEX root=../../main.tex

This chapter combines ideas from the literature on convex neural networks,
accelerated methods, and optimization algorithms for constrained problems.
We now briefly survey these concepts.

\textbf{Convex Reformulations}:
The most similar work in the literature is by \citet{bai2022efficient}, who
consider training two-layer ReLU networks via convex reformulations using ADMM.
Their approach requires solving a linear system at each iteration,
or uses coordinate descent to solve the ADMM sub-problems.
In practice, our solvers scale to larger datasets
and allow for more activation patterns.

In subsequent work, \citet{feng2024chronos} propose Cronos, an ADMM solver for
convex reformulations based on Nystr\"{o}m preconditioning.
Their work emphasizes GPU acceleration and scales convex optimization for
neural networks to ImageNet, but also relies on aggressively sub-sampling the
convex program for efficiency.
Compared to our approach, Cronos trades global optimality guarantees for
enhanced speed and scalability.

Another subsequent work by \citet{sahiner2024burer} explores training ReLU
networks with vector-outputs.
While we use a one-vs-all prediction strategy to avoid the challenging
co-positive program arising in the convex reformulation
\citep{sahiner2021vector}, they rely on Burer-Monteiro factorization to obtain
a benign non-convex program \citep{burer2005low}.
They minimize the resulting non-convex problem with SGD, which is
(theoretically) much slower than our accelerated convex solvers.

\textbf{Accelerated Proximal Gradient}:
\citet{beck2009fista, nesterov2007proximalgradient} were the first to extend
optimal gradient methods~\citep{nesterov1983method} to composite problems.
Work since then includes extensions to
stochastic~\citep{schmidt2011convergence} and
non-convex~\citep{li2015accelerated} optimization.
See~\citet{parikh2014proximalsurvery} for a survey of proximal algorithms,
including proximal gradient.

\textbf{Augmented Lagrangian Methods}:
The convergence theory was initially developed by
Rockafellar~\citep{rockafellar1976augmented, rockafellar1976monotone}.
More recent work includes practical guidelines~\citep{birgin2014pratical} and
acceleration techniques~\citep{kang2015inexact}.
See \citet{bertsekas2014constrained} for exhaustive theoretical developments.

